\section{Method}
\label{sec:method}

Section~\ref{sec:setup} defines the candidate universe, forgetting gate,
realized damage, and sealed information boundary. We now define the two
profile coordinates, their combined score, and the common repair operator.
All target-specific profile quantities are computed at $\theta_0$ and sealed
before the parent trajectory begins.

\subsection{Loss Susceptibility: A Gradient-Based Signal}
\label{sec:method-loss-susceptibility}
\label{sec:method-profile}

\begin{figure}[!t]
\centering
\begingroup
\definecolor{lsViolet}{HTML}{6A51A3}
\definecolor{lsOrange}{HTML}{D95F02}
\definecolor{lsGray}{HTML}{737373}

\resizebox{0.98\columnwidth}{!}{%
\begin{tikzpicture}[
  x=1cm,
  y=1cm,
  line cap=round,
  line join=round,
  >={Latex[length=1.7mm,width=1.05mm]},
  every node/.style={font=\sffamily},
  probevec/.style={
    -{Latex[length=1.55mm,width=.98mm]},
    draw=lsViolet!74,
    line width=.72pt
  },
  pairvec/.style={
    -{Latex[length=1.75mm,width=1.10mm]},
    draw=lsOrange,
    line width=1.05pt
  },
  response/.style={
    draw=lsViolet!40,
    fill=lsViolet!5,
    rounded corners=.65mm,
    minimum width=1.12cm,
    minimum height=.47cm,
    inner sep=.55mm,
    font=\sffamily\fontsize{6.4}{7.0}\selectfont
  },
  result/.style={
    draw=lsOrange!68!black,
    fill=lsOrange!5,
    rounded corners=.8mm,
    minimum width=1.55cm,
    minimum height=1.06cm,
    align=center,
    inner sep=.75mm,
    font=\sffamily\fontsize{6.3}{7.0}\selectfont
  },
  smalllabel/.style={
    text=lsGray,
    font=\sffamily\fontsize{6.1}{6.7}\selectfont
  }
]

\def\R{1.18}
\coordinate (origin) at (1.42,1.67);

% ------------------------------------------------------------------
% One shared fixed-radius neighborhood
\fill[lsViolet!4]
  (origin) circle (\R);

\draw[lsGray!58,line width=.70pt]
  (origin) circle (\R);

% ------------------------------------------------------------------
% Multiple paired random directions on the same sphere
\foreach \ang in {7,51,93,136,177}{
  \draw[probevec]
    (origin)
    --
    ({1.42+\R*cos(\ang)},
     {1.67+\R*sin(\ang)});

  \draw[probevec]
    (origin)
    --
    ({1.42-\R*cos(\ang)},
     {1.67-\R*sin(\ang)});

  \fill[lsViolet!78]
    ({1.42+\R*cos(\ang)},
     {1.67+\R*sin(\ang)})
    circle[radius=.030];

  \fill[lsViolet!78]
    ({1.42-\R*cos(\ang)},
     {1.67-\R*sin(\ang)})
    circle[radius=.030];
}

% ------------------------------------------------------------------
% Highlight one symmetric probe pair
\def\pairang{28}

\draw[pairvec]
  (origin)
  --
  ({1.42+\R*cos(\pairang)},
   {1.67+\R*sin(\pairang)});

\draw[pairvec]
  (origin)
  --
  ({1.42-\R*cos(\pairang)},
   {1.67-\R*sin(\pairang)});

\fill[lsOrange]
  ({1.42+\R*cos(\pairang)},
   {1.67+\R*sin(\pairang)})
  circle[radius=.042];

\fill[lsOrange]
  ({1.42-\R*cos(\pairang)},
   {1.67-\R*sin(\pairang)})
  circle[radius=.042];

% Frozen checkpoint
\fill[black!82]
  (origin) circle[radius=.052];

\node[
  anchor=north,
  text=black!82,
  font=\sffamily\bfseries\fontsize{6.6}{7.2}\selectfont
] at (1.42,1.59)
  {$\theta_0$};

% Symmetric endpoints
\node[
  anchor=south west,
  text=lsOrange,
  font=\sffamily\bfseries\fontsize{5.9}{6.5}\selectfont
] at (
  {1.42+\R*cos(\pairang)+.03},
  {1.67+\R*sin(\pairang)+.02}
)
  {$+\eta v_r$};

\node[
  anchor=north east,
  text=lsOrange,
  font=\sffamily\bfseries\fontsize{5.9}{6.5}\selectfont
] at (
  {1.42-\R*cos(\pairang)-.03},
  {1.67-\R*sin(\pairang)-.02}
)
  {$-\eta v_r$};

\node[
  anchor=north,
  align=center,
  smalllabel
] at (1.42,.30)
  {one shared sphere\\
   $\|\Delta\theta_{\mathcal B}\|_2=\eta$};

% ------------------------------------------------------------------
% Direction-specific squared responses
\node[response] (d1) at (3.78,2.57)
  {$\delta_1(x)^2$};

\node[response] (d2) at (3.78,1.97)
  {$\delta_2(x)^2$};

\node[
  text=lsGray,
  font=\sffamily\fontsize{8.0}{8.5}\selectfont
] at (3.78,1.39)
  {$\vdots$};

\node[response] (dR) at (3.78,.82)
  {$\delta_R(x)^2$};

\draw[
  -{Latex[length=1.55mm,width=.98mm]},
  lsGray!75,
  line width=.65pt
] (2.75,1.67) -- (3.12,1.67);

% ------------------------------------------------------------------
% Aggregate over all R directions
\node[result] (susceptibility) at (5.88,1.67)
  {\textbf{Loss Susceptibility}\\[.25mm]
   $\displaystyle
    \widehat q_G(x)
    =
    \frac{d_{\mathcal B}}{R}
    \sum_{r=1}^{R}\delta_r(x)^2$};

\draw[
  -{Latex[length=1.65mm,width=1.02mm]},
  lsOrange!78!black,
  line width=.78pt
] (4.42,1.67) -- (susceptibility.west);

\path[use as bounding box]
  (0.05,.02) rectangle (6.90,3.08);

\end{tikzpicture}%
}
\endgroup

\caption{\textbf{Fixed-radius Loss Susceptibility probes.}
All $R$ random block directions lie on one radius-$\eta$ neighborhood around
$\theta_0$. Each direction is evaluated symmetrically at
$\theta_0\pm\eta v_r$. The dimension-scaled mean of the squared
central-difference responses estimates Loss Susceptibility.}
\Description{A single fixed-radius sphere centered at the frozen checkpoint
contains multiple paired random directions. Each pair produces a squared
central-difference directional response, and their dimension-scaled mean
estimates Loss Susceptibility.}
\label{fig:loss-susceptibility-geometry}
\end{figure}

The realized target update is unavailable at profiling time and may change
direction over a multi-step trajectory. Without conditioning on that future
direction, we define the \emph{Block-Local Loss Susceptibility} of retained
behavior $x$ for a declared parameter block $\mathcal B$ as
\begin{equation}
  \begin{aligned}
  q_G^\star(x)
  &=\|\nabla_{\mathcal B}\ell_x(\theta_0)\|_2^2 \\
  &=\sup_{\substack{\|u\|_2=1\\
       \operatorname{supp}(u)\subseteq\mathcal B}}
       |D_u\ell_x(\theta_0)|^2 .
  \end{aligned}
  \label{eq:block-local-loss-susceptibility}
\end{equation}
Thus $q_G^\star$ is the squared $\ell_2$-norm of the loss gradient restricted
to $\mathcal B$, equivalently the largest squared first-order response over
unit-norm block perturbations. It depends on the parameterization, norm, and
chosen block, so we compare rankings within a setting rather than magnitudes
across models. The exact quantity is therefore gradient-defined; only its
estimator $\widehat q_G$ is forward-only.

Let $d_{\mathcal B}$ be the block dimension. Draw
$g_r\overset{\mathrm{iid}}{\sim}\mathcal N(0,I_{d_{\mathcal B}})$, normalize
$v_r=g_r/\|g_r\|_2$, and set it to zero outside $\mathcal B$. For radius
$\eta>0$ and $R\in\mathbb N$ directions, define the forward-only estimate
\begin{equation}
\begin{aligned}
  \delta_r(x)
  &=
  \frac{\ell_x(\theta_0+\eta v_r)-\ell_x(\theta_0-\eta v_r)}
       {2\eta},\\
  \widehat q_G(x)
  &=
  \frac{d_{\mathcal B}}{R}\sum_{r=1}^{R}\delta_r(x)^2,
  \qquad
  \mathbb E_{v_{1:R}}[\widehat q_G(x)]
  =q_G^\star(x)+O(\eta^2).
\end{aligned}
\label{eq:loss-susceptibility-estimator}
\end{equation}
The leading-order relation follows from the symmetric directional derivative
and $\mathbb E[v_rv_r^\top]=I/d_{\mathcal B}$, with bounded third directional
derivatives near $\theta_0$; the expectation holds over the probe bank with
$x$, $\theta_0$, and $\mathcal B$ fixed. In the directional-derivative limit,
the probe-bank variance is
$2(d_{\mathcal B}-1)\bigl(q_G^\star(x)\bigr)^2/[R(d_{\mathcal B}+2)]$.
Thus $\eta$ controls central-difference bias and finite $R$ controls
Monte Carlo variability. Figure~\ref{fig:loss-susceptibility-geometry}
illustrates that all responses are sampled from the same fixed-radius
neighborhood.

Each direction requires two batched forward sweeps over the candidate pool,
starting from and returning to the same parameter copy. Target profiling
therefore uses $2R$ pool-forward sweeps, materializes no per-candidate
gradient, and requires no candidate-side backward pass. A
target-request-disjoint calibration fold selects $(R,\eta)$ under the
predeclared rule using rank agreement and high-rank overlap with $q_G^\star$,
plus perturbation-survival and split-bank checks. Direct gradients
are streamed only on the declared fidelity support. The resulting rank-fidelity
certificate against the direct-gradient reference is setting-level; it does
not certify each target candidate support. Each target profile instead must
pass perturbation-survival, split-bank rank-agreement, and integrity checks on
its discovery fold before
the sealed audit is opened. This provides a backward-free target access pattern
that never materializes candidate-specific gradients, not a universal claim
of fewer FLOPs or lower wall-clock time.

\subsection{Representation Proximity and the Sealed Score}
\label{sec:method-representation}
\label{sec:method-mixture}
\label{sec:method-channels}

Representation Proximity measures exposure to the current forget request in
hidden-state space. Let $h(x;\theta_0)$ be the mean final hidden state over
answer tokens only, and let $\mathcal N_k(x;\mathcal D_f)$ contain the $k$
forget examples with highest cosine similarity to $x$. We define
\begin{equation}
  \widehat q_H(x)
  =
  \frac{1}{k}\sum_{z\in\mathcal N_k(x;\mathcal D_f)}
  \operatorname{cos}\!\left(h(x;\theta_0),h(z;\theta_0)\right).
  \label{eq:representation-proximity}
\end{equation}
This forward-only coordinate is a proxy for request exposure, not a claim of
causal dependence.

Let $N_{\mathrm{disc}}=|\mathcal C_{\mathrm{disc}}|$. For
$j\in\{G,H\}$, empirical midrank maps fitted
on $\mathcal C_{\mathrm{disc}}$ put the coordinates on a common $[0,1]$ scale:
\begin{equation}
\begin{aligned}
F_j^{\mathrm{disc}}(z)
&=\frac{1}{N_{\mathrm{disc}}}\sum_{y\in\mathcal C_{\mathrm{disc}}}
\left(\mathbf 1\{\widehat q_j(y)<z\}+\tfrac12\mathbf 1\{\widehat q_j(y)=z\}\right),\\
\widetilde q_j(x)&=F_j^{\mathrm{disc}}(\widehat q_j(x)).
\end{aligned}
\label{eq:midrank-normalization}
\end{equation}
Their mixture is
\begin{equation}
  S_\alpha(x)
  =(1-\alpha)\widetilde q_G(x)+\alpha\widetilde q_H(x),
  \qquad \alpha\in[0,1],
  \label{eq:mixture}
\end{equation}
with $S_0=\widetilde q_G$ denoting the estimated Loss Susceptibility endpoint
obtained by rank-normalizing $\widehat q_G$, and $S_1=\widetilde q_H$ denoting the
Representation Proximity endpoint.

For each setting--parent cell, a target-request-disjoint development fold
selects $\widehat\alpha_{\mathrm{pred}}$ by average within-request rank
correlation, with the tie rule fixed in advance. This is scalar calibration
from completed development trajectories, not outcome-free learning. The
target score $S_{\widehat\alpha_{\mathrm{pred}}}$ then fixes both the audit
ranking and the Top-$K_p$ eligible discovery pool before target unlearning.
Audit outcomes select neither the weight nor the pool.

\subsection{Protection under a Fixed Selection Quota}
\label{sec:method-protection}

Protection asks whether the sealed score can allocate a limited retain-side
repair budget. Every selector chooses exactly $K_p$ eligible examples from
$\mathcal C_{\mathrm{disc}}$. The processed-token limit
$B_{\mathrm{tok}}$ separately fixes repair compute. Starting from the common
entry checkpoint $\theta_{t^\dagger}$, the repair objective for selected pool
$\mathcal P$ is
\begin{equation}
  \begin{aligned}
    L_{\mathrm{rep}}(\theta)
    &=
    \frac{1}{|\mathcal P|}
    \sum_{x\in\mathcal P}\ell_x(\theta)
    +\beta L_{\mathrm{anc}}(\theta),\\
    L_{\mathrm{anc}}(\theta)
    &=
    \frac{1}{|\mathcal R_0|}
    \sum_{x\in\mathcal R_0}
    \operatorname{KL}\!\left(
      p_{\theta_{t^\dagger}}(\cdot|x)\,\|\,p_\theta(\cdot|x)
    \right),
    \qquad \beta\geq0 ,
  \end{aligned}
  \label{eq:budgeted-repair}
\end{equation}
where the selector-independent neutral stream $\mathcal R_0$ anchors the
entry checkpoint.

All active arms share the operator, neutral stream, example order, optimizer
random stream, token budget, save schedule, rollback rule, and
forgetting/utility guards. Only the selected pool changes. The returned model
is the last saved checkpoint satisfying the predeclared feasibility region;
otherwise the arm is infeasible. The two-signal selector is compared with
$S_0$, $S_1$, and repeated random selection under this common harness. The
unrepaired entry checkpoint separately establishes whether repair has net
benefit. Confirmatory protection is measured on
$\mathcal C_{\mathrm{audit}}$, never on the score-selected pool. The full guard
implementation is given in Appendix~\ref{app:repair-protocol}.
